\chapter{Six-Robot Cluster Kinematics}
\label{app:kin6}

% New text. Checked against src/control/pentagon_kinematics.py
% (VF_Robot), which implements this parameterization.
The six-robot cluster of Chapter~\ref{ch:second_order} is a cluster of
clusters. The robots form three pairs, and the pair midpoints form a
triangle described with the same SAS variables as the three-robot cluster
of Appendix~\ref{app:kin3}. The estimator of
Chapter~\ref{ch:second_order} reads only the six measured positions, so
this parameterization matters to the formation controller alone.

\section{Forward Kinematics}
\label{sec:appC:forward}

Pair A holds robots 1 and 2, pair B robots 3 and 4, and pair C robots 5
and 6. The pair midpoints are
\begin{equation}
    \mathbf{m}_A = \tfrac{1}{2}(\mathbf{p}_1 + \mathbf{p}_2), \quad
    \mathbf{m}_B = \tfrac{1}{2}(\mathbf{p}_3 + \mathbf{p}_4), \quad
    \mathbf{m}_C = \tfrac{1}{2}(\mathbf{p}_5 + \mathbf{p}_6),
    \label{eq:appC:midpoints}
\end{equation}
and the cluster centroid is their mean,
$\mathbf{p}_c = \tfrac{1}{3}(\mathbf{m}_A + \mathbf{m}_B + \mathbf{m}_C)$.
The midpoint triangle has sides
$p_1 = \lVert\mathbf{m}_B - \mathbf{m}_A\rVert$ and
$q_1 = \lVert\mathbf{m}_C - \mathbf{m}_B\rVert$, with $\beta_1$ the
interior angle at $\mathbf{m}_B$, and the heading $\theta_c$ is the
direction from $\mathbf{m}_A$ to $\mathbf{m}_B$. Each pair has a length
and an orientation,
\begin{equation}
    L_2 = \lVert\mathbf{p}_2 - \mathbf{p}_1\rVert, \qquad
    \theta_2 = \operatorname{atan2}(y_2 - y_1,\, x_2 - x_1),
    \label{eq:appC:pair}
\end{equation}
and likewise $(L_3, \theta_3)$ for pair B and $(L_4, \theta_4)$ for pair C.
The twelve cluster variables are
\begin{equation}
    (x_c,\, y_c,\, \theta_c,\, p_1,\, \beta_1,\, q_1,\,
     L_2,\, \theta_2,\, L_3,\, \theta_3,\, L_4,\, \theta_4).
    \label{eq:appC:state}
\end{equation}

\section{Inverse Kinematics}
\label{sec:appC:inverse}

The inverse kinematics builds the midpoint triangle from
$(p_1, \beta_1, q_1)$ in a local frame, as in
Appendix~\ref{app:kin3}, then rotates it by the heading and translates it
to $(x_c, y_c)$. Each pair's two robots are then placed at its midpoint
plus and minus $\tfrac{1}{2}L_k(\cos\theta_k, \sin\theta_k)$.

\section{Inverse Jacobian}
\label{sec:appC:jacobian}

The $12 \times 12$ inverse Jacobian $\mathbf{J}_c^{-1}$ maps the rates of
the twelve cluster variables to the velocities of the six robots. Its
entries are the partial derivatives of the inverse kinematics and are not
listed here. The simulation software computes them analytically
(\texttt{compute\_inverse\_jacobian} in
\texttt{src/control/pentagon\_kinematics.py}), and the
\texttt{cluster\_builder} tool generates them in closed symbolic form
(\texttt{clusterbuilder.py} with the \texttt{--symbolic} option). The
cluster's heading is not regulated, since the formation controller
commands zero spin throughout.
